% 本片段自印刷页 p.79 顶部接续上一片段末句（"……会以有趣的方式简化"），原文在此直接给出编号公式 (2.335)。

\begin{equation}
  Z(z,\tau;\bar z,\bar\tau)
  = e^{-4\pi k\,\frac{(\operatorname{Im}z)^2}{\operatorname{Im}\tau}}
  \sum_{r,\tilde r\in\mathbb{Z}/2k\mathbb{Z}}
  D_{r,\tilde r}\,\chi_r(z,\tau)\,\overline{\chi_{\tilde r}(z,\tau)}
  \tag{2.335}
  \label{eq:2.335}
\end{equation}
% [Tutor Note: (2.335) 的前因子指数是 $(\operatorname{Im}z)^2/\operatorname{Im}\tau$，不是 $|z|^2$ 也不是 $z\bar z$：它来自环面上与 $z$ 相关的零模方向的高斯积分，故与 $\tau$ 的虚部绑定。上划线只覆盖右移因子 $\chi_{\tilde r}$，作者并没有把它写成 $|\chi_r|^2$，因为左移指标 $r$ 与右移指标 $\tilde r$ 一般并不相等；把 (2.335) 与 (2.328) 并排对照时，最容易错的就是这一点。]

矩阵 $D_{r,\tilde r}$ 将在下文第 2.4.1 节中进一步讨论。

由上述计算我们可以得出若干重要结论：

\begin{enumerate}
  \item 在 \eqref{eq:2.328} 与 \eqref{eq:2.335} 两式中，我们都看到共形场论的典型特征：配分函数是若干「全纯对象乘以反全纯对象」之和，至多差一个简单的局部因子。注意它们对 $z$ 与 $\tau$ 都是全纯的。这些全纯因子称为\emph{共形块}（conformal blocks），而\emph{有理}共形场论的一条（定义性的）性质是：所有黎曼面上的所有关联函数都是这样的和。该局部前因子代表了在同时要求规范不变性与全纯分解时所出现的\emph{反常}（anomaly）。这一观察至少可以追溯到 [5]。

  {\small\itshape ［原文标记：确实需要讨论模变换，才能看出真正不变的东西只是共形块构成的向量空间。］}
  % [原文待办标记，原文 p.79；位于页边竖排]

  \item 再者，作为 $z$ 的函数，共形块的向量空间与 $U(1)_k$ Chern-Simons 理论的 Hilbert 空间相同。这将成为一个普遍原理：

  \begin{quote}
    \itshape
    在 $\Sigma\times\mathbb{R}$ 上量子化的 Chern-Simons 规范理论中，物理态的向量空间，与相应的有理共形场论的共形块空间之间存在同构。此外，若该 Chern-Simons 理论是借助由 $\Sigma$ 上某个复结构的选择所诱导的 Kähler 量子化来量子化的，则这乃是黎曼面模空间上全纯向量丛之间的等同。
  \end{quote}

  我们的目标之一，就是在更一般的 Chern-Simons 理论的背景下更好地理解这一陈述。

  \item 在这个意义上，带边的三维流形上的拓扑 Chern-Simons 理论与边界上的（有理共形场论的手征一半）之间存在某种等价。这是一种「全息」（holography）。
\end{enumerate}

\subsubsection{把同一手征代数粘合出不同半径} % 原文: 2.4.1 Gluing The Same Chiral Algebra To Get Different Radii

爱钻研的读者会注意到，整数 $k$ 可以有多种分解 $k=pq$（即使 $k$ 是素数时也是如此）。那么是什么挑出了不同的半径？

\underline{在 RCFT 框架内的回答}

在 \eqref{eq:2.335} 中，矩阵 $D_{r,r'}$ 告诉我们如何把左移与右移部分粘合起来。如 [43] 所示，这必须通过 $\mathbb{Z}/(2k\mathbb{Z})$ 融合规则代数的一个自同构来完成。（这是关于 RCFT 中表示粘合的一般性结果。）把该阿贝尔群按加法来想，自同构就是 $x\mapsto yx$，其中 $y$ 与 $2k$ 互素。[43] 中声称：若取基本区 $0\le r\le 2k-1$，则当 $r=pr_2+qr_1$、$r'=qr_1-pr_2$，而 $(r_1,r_2)$ 取值于 $\mathbb{Z}\oplus\mathbb{Z}$ 中由 $(\pm p,0)$ 与 $(0,\pm q)$ 所决定的菱形区域时，$D_{r,r'}=1$。

{\small\itshape ［原文标记：最好再核对一遍。有一篇 Dijkgraaf–Verlinde 的相关论文，需要引用。］}
% [原文待办标记，原文 p.80；位于页边竖排]
% [Tutor Note: 作者自己在页边写了「最好再核对一遍」，说明 $D_{r,r'}$ 的这条「菱形区域」判据在原稿中尚未最终定稿。它的含义是把融合规则代数的自同构 $x\mapsto yx$ 图解化：$r\mapsto r'$ 必须是该自同构诱导的置换，而菱形区域只是保证 $(r_1,r_2)\mapsto(pr_2+qr_1,\,qr_1-pr_2)$ 在其中单值。引用 $D_{r,r'}$ 的具体形式（例如计算某个 $k=pq$ 下的显式矩阵）之前，务必先用 $\gcd(y,2k)=1$ 核对它确实是双射。]

\underline{在 Maxwell-Chern-Simons 理论框架内的回答}

正如 [26] 中所详加讨论的，在 Chern-Simons 理论的框架内导出最低 Landau 能级（lowest Landau level, LLL）波函数，实际上会挑出一个优先的半径。让我们回到 \eqref{eq:2.199}，把 LLL 中波函数的公式代入。我们得到：
\begin{equation}
  \bar\Psi(A) := \mathcal{N}\,e^{-4\pi k\operatorname{Im}\tau\,A_zA_{\bar z}}
  \sum_{\omega\in\mathcal{H}^1_{\mathbb{Z}}}
  e^{-4\pi k\operatorname{Im}\tau\,\omega_z\omega_{\bar z}}
  e^{-8\pi k\operatorname{Im}\tau\,\omega_{\bar z}A_z}
  \,\psi(A_z+\omega_z)
  \tag{2.336}
  \label{eq:2.336}
\end{equation}
% [Tutor Note: 求和指标 $\omega$ 属于 $\mathcal{H}^1_{\mathbb{Z}}$，即整系数上同调类所标记的大平移/大动量扇区。指数里 $\omega_z\omega_{\bar z}$ 与 $\omega_{\bar z}A_z$ 的系数分别是 $4\pi k$ 与 $8\pi k$，比值 $1:2$ 正是把二次型在 $A+\omega$ 处展开后交叉项的系数；把 (2.336) 与 (2.199) 对照或做配方时若丢掉这个 $1:2$，就得不到正确的半径。]

我们能认出来，这是某个周期标量场（同时含左移与右移自由度）在某个\underline{特定半径}处的瞬子求和，即
\begin{equation}
  S=\frac{k}{4\pi}\int d\phi*d\phi
  \tag{2.337}
  \label{eq:2.337}
\end{equation}
其中 $\phi\sim\phi+2\pi$，并与 $A^{1,0}$ 有手征耦合：
\begin{equation}
  2ik\int\bar\partial\phi\wedge A^{1,0}
  \tag{2.338}
  \label{eq:2.338}
\end{equation}

我们可以通过把环面上该玻色子（作用量为上式）的配分函数与 \eqref{eq:2.336} 对照，来解读这一信息。半径为 $r^2=k\ell_s^2$。进一步的讨论（包括向环面规范群的非凡推广）见 [26]。

\underline{在拓扑场论框架内的回答：面缺陷}

融合规则的非平凡自同构，可以解释为在三维理论中定义\emph{面缺陷}（surface defects）。见 [35]。

{\small\itshape ［原文标记（星号框）：关于 Kapustin–Saulina 的图景，需要再多解释一些。］}
% [原文待办标记，原文 p.80]

\subsubsection{自由费米子半径：能级 1/2 的 theta 函数与环面上的旋量结构} % 原文: 2.4.2 The Free Fermion Radius: Level 1/2 Theta Functions And Spin Structures On The Torus

从 \eqref{eq:2.310} 可以清楚地看到，有一些特别值得关注的半径：

\begin{enumerate}
  \item 有著名的 $r\to1/r$ T-对偶对称性，因此我们只需考察 $r\ge1$。

  \item 在 $T$-对偶的不动点 $r=1$ 处，存在增强的 $SO(4)=(SU(2)_L\times SU(2)_R)/\mathbb{Z}_2$ 对称性。

  \item 在半径 $r^2=2$ 处，有 $\sqrt{2pq}=2$，于是我们有共形维数为 $1/8$ 的全纯算符 $e^{\pm\frac{i}{2}X_L(z)}$，以及它们的平方
  \begin{equation}
    \psi(z)=e^{iX_L(z)}
    \qquad
    \bar\psi(z)=e^{-iX_L(z)}
    \tag{2.339}
    \label{eq:2.339}
  \end{equation}
  其维数为 $h=1/2$。

  \item 在半径 $k=6$ 处，理论实际上具有 $N=2$ 超对称。$r=\pm6$ 的表示给出
  \begin{equation}
    G^{\pm}(z)\sim e^{\pm i\sqrt{3}X}
    \tag{2.340}
    \label{eq:2.340}
  \end{equation}
  它们惊人地满足 $N=2$ 超共形代数。
\end{enumerate}

最后两个例子启发了引入\emph{超手征代数}的想法。在自由费米子情形，我们有
\begin{equation}
  \begin{aligned}
    \psi(z_1)\psi(z_2) &\sim (z_1-z_2)\,{:}\,\psi(z_1)\psi(z_2)\,{:}\\
    \psi(z_1)\bar\psi(z_2) &\sim (z_1-z_2)^{-1}+{:}\,\partial\psi\,\psi\,{:}+\mathcal{O}(z_1-z_2)
  \end{aligned}
  \tag{2.341}
  \label{eq:2.341}
\end{equation}

于是 $\psi$ 与 $\bar\psi$「几乎」构成一个手征代数：它们的关联函数是单值的，但存在非平凡的编织：关联函数中 $\psi$ 与 $\bar\psi$ 的排序是有影响的。

\begin{enumerate}[start=3]
  \item 无论如何，在这个特殊的「自由费米子半径」处，高斯模型存在一个「双重覆盖」。然而，相当显然地，要定义关联函数我们必须定义旋量结构。

  \item 这或许最明显的例子就在配分函数的计算中。首先，我们必须为
  \begin{equation}
    \psi(z)=\sum d_rz^{-r-1/2}dz^{1/2}
    \tag{2.342}
    \label{eq:2.342}
  \end{equation}
  选取费米子边界条件，其中 $z=e^{i\phi+t}$，视作 $\phi$ 的函数。
  % [Tutor Note: (2.342) 的求和号下原文本就未写出指标范围：NS 扇区取 $r\in\mathbb{Z}+\frac12$，R 扇区取 $r\in\mathbb{Z}$（见本节下一条与附录 G）。此外因子 $dz^{1/2}$ 表明 $\psi$ 是 $\frac12$-微分；在 $z=e^{i\phi+t}$ 这类坐标变换下，半整数权会带来额外的相位，讨论环面边界条件时不能漏掉它。]

  \item 如附录 G 中所解释的，圆上有两种旋量结构，即 Neveu-Schwarz（NS）与 Ramond（R）旋量结构，分别对应于 $\psi$ 的反周期与周期边界条件。用模式展开来说，NS 对应于 $r\in\mathbb{Z}+\frac12$，R 对应于 $r\in\mathbb{Z}$。

  \item $\operatorname{Tr}e^{-\beta H}$ 自然地选取反周期边界条件。[见任何一本 QFT 教科书，或自行验证]。

  \item 总共有四个自然的环面配分函数，对应于环面上的四种旋量结构。（约定：$Z(g_t,g_s;\tau)$ 表示环面上带扭曲边界条件。）

  {\small\itshape ［原文标记：弄个方框宏（box macro）来画带边界条件的环面。］}
  % [原文待办标记，原文 p.81；位于页边竖排]

  \begin{equation}
    \begin{aligned}
      Z(+,+) &= \frac{\vartheta_1(z,\tau)}{\eta(\tau)}\\
      Z(-,+) &= \frac{\vartheta_2(z,\tau)}{\eta(\tau)}\\
      Z(-,-) &= \frac{\vartheta_3(z,\tau)}{\eta(\tau)}\\
      Z(+,-) &= \frac{\vartheta_4(z,\tau)}{\eta(\tau)}
    \end{aligned}
    \tag{2.343}
    \label{eq:2.343}
  \end{equation}

  {\small\itshape ［原文标记：ETC.（等等）］}
  % [原文标记，原文 p.82 首行；作者此处以大写 ETC. 表示「以下同理」，属未写完的列表收尾]

  这里我们引入了半整数能级的 theta 函数。

  定义：

  {\small\itshape ［原文标记：取决于我们在「带电粒子绕环运动」那一节里讲了多少，此处就直接回指那里。］}
  % [原文待办标记，原文 p.82；位于页边竖排]

  \begin{equation}
    \vartheta\left[\genfrac{}{}{0pt}{}{\theta}{\phi}\right](z|\tau)
    := \sum_{n\in\mathbb{Z}}
    e^{i\pi\tau(n+\theta)^2+2\pi i(n+\theta)(z+\phi)}
    \tag{2.344}
    \label{eq:2.344}
  \end{equation}

  \begin{equation}
    \begin{aligned}
      \vartheta_1(z|\tau) &:= \vartheta\left[\genfrac{}{}{0pt}{}{1/2}{1/2}\right](z|\tau)
      = 2q^{1/8}\cos\left[\pi\left(z+\tfrac12\right)\right]+2q^{9/8}\cos\left[3\pi\left(z+\tfrac12\right)\right]+\cdots\\
      \vartheta_2(z|\tau) &:= \vartheta\left[\genfrac{}{}{0pt}{}{1/2}{0}\right](z|\tau)
      = 2q^{1/8}\cos(\pi z)+2q^{9/8}\cos(3\pi z)+\cdots\\
      \vartheta_3(z|\tau) &:= \vartheta\left[\genfrac{}{}{0pt}{}{0}{0}\right](z|\tau)
      = 1+2\cos[2\pi z]\,q^{1/2}+2\cos[4\pi z]\,q^{2}+\cdots\\
      \vartheta_4(z|\tau) &:= \vartheta\left[\genfrac{}{}{0pt}{}{0}{1/2}\right](z|\tau)
      = 1-2q^{1/2}\cos[2\pi z]+2q^{2}\cos[4\pi z]+\cdots
    \end{aligned}
    \tag{2.345}
    \label{eq:2.345}
  \end{equation}
  % [Tutor Note: (2.345) 用的是 Jacobi 的经典约定：$q=e^{2\pi i\tau}$，特征 $\vartheta\bigl[\genfrac{}{}{0pt}{}{\theta}{\phi}\bigr]$ 的上下指标与 $\vartheta_{1},\dots,\vartheta_{4}$ 的对应依次为 $(1/2,1/2)$、$(1/2,0)$、$(0,0)$、$(0,1/2)$。注意 $\vartheta_1$ 的宗量是 $\cos[\pi(z+\frac12)]$ 而 $\vartheta_2$ 是 $\cos(\pi z)$，二者相差一个 $z\to z+\frac12$ 的平移；抄写时最易在此处丢掉 $\frac12$。]

  {\small\itshape ［原文标记：还要给出无穷乘积表示。这些公式都应当放进附录。］}
  % [原文待办标记，原文 p.82；位于页边竖排]

  \item 讨论 $PSL(2,\mathbb{Z})$ 的三个指数为 3 的子群，它们各自保持三种偶旋量结构。用它来解释上述 theta 函数的 $S$ 变换与 $T$ 变换。

  \item 评述自旋算符及其关联函数。
\end{enumerate}

\subsection{旋量理论} % 原文: 2.5 Spin Theories

Chern-Simons 理论的一个重要变体是\emph{旋量 Chern-Simons}（spin-Chern-Simons）理论。在此情形中，$M_3$ 除了定向之外还具有一个\emph{旋量结构}，而 Chern-Simons 作用量确实依赖于该旋量结构。

旋量结构的数学背景见附录 G。

我们可以借助把规范丛配边到零的办法，来定义闭三维流形上的 Chern-Simons 不变量。这意味着我们要找到一个四维流形 $W_4$，使得
\begin{equation}
  \partial W_4=M_3
  \tag{2.346}
  \label{eq:2.346}
\end{equation}
并把丛延拓到 $W_4$ 上。丛一旦可以延拓，由单位分解的论证，延拓联络就没有任何问题。现在，有了这些额外数据，我们就可以把 Chern-Simons 不变量
\begin{equation}
  CS(A)=\frac{1}{(2\pi)}\int_{M_3}AdA\ \in\mathbb{R}/2\pi\mathbb{Z}
  \tag{2.347}
  \label{eq:2.347}
\end{equation}
提升为真正的实数：
\begin{equation}
  \widehat{CS}(A;W_4)=2\pi\int_{W_4}\frac{F}{2\pi}\wedge\frac{F}{2\pi}\ \in\mathbb{R}
  \tag{2.348}
  \label{eq:2.348}
\end{equation}
% [Tutor Note: (2.347) 用的是阿贝尔情形 $U(1)$ 的归一化 $\frac{1}{2\pi}\int A\,dA$（既无 $\operatorname{tr}$，也不是 $\frac{1}{4\pi}$），与非阿贝尔的 (3.7)（$\frac{k}{8\pi^2}\operatorname{tr}$）是两套完全不同的约定；把 (2.354) 的相位因子与 (3.12) 的 Wess-Zumino 型整数对比之前，必须先统一到同一套归一化，否则会莫名其妙地差出因子 2 或 4。]

$W_4$ 的选取可以不同，而这些差别
\begin{equation}
  \left(\widehat{CS}(A;W_4)-\widehat{CS}(A;W'_4)\right)
  =2\pi\int_{Z_4}\frac{F}{2\pi}\wedge\frac{F}{2\pi}\ \in2\pi\mathbb{Z}
  \tag{2.349}
  \label{eq:2.349}
\end{equation}
% 原 (2.349) 把第二个四维流形印作 $W'_r$（下标为 $r$ 而非 $4$），紧随其后的正文又称之为 $W'_4$，这是原稿的排版笔误；此处照原样转写，引用时应按 $W'_4$ 理解。
是某个闭四维流形 $Z_4$（由把 $W_4$ 与 $W'_r$ 粘合（gluging）而得）上，丛的第一 Chern 类的代表元之楔积的积分。这是一个整形式的积分，因而是一个整数。例如，若 $Z_4=\mathbb{CP}^2$ 且 $[F/2\pi]=c_1(H)$，其中 $H$ 是基本的超平面丛（tautological hyperplane bundle），则该整数为 $+1$。

{\small\itshape ［原文标记：评注一下对 $A$ 延拓到 $W_4$ 的方式的依赖。］}
% [原文待办标记，原文 p.82；位于页边竖排]

我们想把这个 CS 作用量除以 2，以定义半整数能级的 Chern-Simons 理论。但这将导致一个定义不良的作用量。上面 $Z_4=\mathbb{CP}^2$ 的例子表明情况确实可能如此。

这「只不过」是一个符号问题，因为我们「只不过」是在取 $\exp[iCS(A)]$ 的平方根；但此处若马虎，会导致严重错误。

$H^2(Z_4;\mathbb{Z})$ 是一个整格。一般地，给定整格 $\Lambda$，满足
\begin{equation}
  v^2=c\cdot v\mod 2
  \tag{2.350}
  \label{eq:2.350}
\end{equation}
的向量 $c$ 称为\emph{特征向量}（characteristic vector）。特征向量总是有无穷多个，而任意两个（比如 $c,c'$）相差一个可被 2 整除的向量，即 $c-c'\in2\Lambda$。当 $c$ 是特征向量时，$\frac12v\cdot(v+c)$ 是整数。

结果表明，$w_2\in H^2(Z_4;\mathbb{Z}_2)$ 的任意整提升 $\widehat{w_2}$ 都是特征向量。这促使我们给出半整数 Chern-Simons 项的定义：

\emph{我们在 $W_4$ 上选取一个旋量结构 $\mathsf{s}$，并令}
\begin{equation}
  \frac12\widehat{CS}(A;W_4;\mathsf{s})
  :=2\pi\int_{W_4}\frac12\frac{F}{2\pi}\wedge\left(\frac{F}{2\pi}+\widehat{w_2}\right)\ \in\mathbb{R}
  \tag{2.351}
  \label{eq:2.351}
\end{equation}

{\small\itshape ［原文标记：你这里用的是 $\widehat{w_2}$。需要说明，同一旋量结构的不同提升无关紧要。］}
% [原文待办标记，原文 p.83；位于页边竖排]
% [Tutor Note: 这里正是作者留待补充的一步。若 $\widehat{w_2}$ 与 $\widehat{w_2}'$ 是 $w_2$ 的两个整提升，则 $\widehat{w_2}-\widehat{w_2}'=2\beta$，其中 $\beta$ 是 $W_4$ 上的整形式；代回 (2.352) 后多出的项为 $2\pi\int_{Z_4}\frac{F}{2\pi}\wedge\beta\in2\pi\mathbb{Z}$，故指数化作用量不变，定义与提升的选取无关。原文把这步留作 "Need to explain"，读者须自行补上。]

现在，若有两个延拓，使得所有数据（包括旋量结构）在公共边界上一致，则
\begin{equation}
  \frac12\widehat{CS}(A;W_4;\mathsf{s})-\frac12\widehat{CS}(A;W'_4;\mathsf{s}')
  =2\pi\int_{Z_4}\frac12\frac{F}{2\pi}\wedge\left(\frac{F}{2\pi}+\widehat{w_2}\right)\ \in2\pi\mathbb{Z}
  \tag{2.352}
  \label{eq:2.352}
\end{equation}
正如所愿。因此我们可以把旋量 Chern-Simons 作用量定义为
\begin{equation}
  \exp\left[2\pi i\frac12CS(A;\mathsf{s})\right]
  :=\exp\left[2\pi\int_{W_4}\frac12\frac{F}{2\pi}\wedge\left(\frac{F}{2\pi}+\widehat{w_2}\right)\right]
  \tag{2.353}
  \label{eq:2.353}
\end{equation}

我们付出的代价是：

\begin{enumerate}
  \item 须把注意力限制在带旋量结构选择的 $W_4$ 上。这不是严重的限制，因为三维的旋量配边群是平凡的，带主 $U(1)$ 丛的旋量流形的配边群亦然。（在更高维，这类限制可能变得相当重要。）

  \item 我们需要把 $w_2$ 的一个显式提升 $\widehat{w_2}$（到 $W_4$ 上的闭 2-形式）作为额外数据纳入进来。

  \item 在边界 $M_3=\partial W_4$ 的邻域内，由于 $M_3$ 容许旋量结构，必有 $w_2(M_3)$ 可平凡化。因此在该邻域内我们可以把 $\widehat{w_2}=d\eta$ 平凡化。然而 $\eta$ 在边界 $\partial W_4$ 上并不为零。$\eta$ 的取值依赖于 $M_3$ 上旋量结构 $\mathsf{s}$ 的选取。两个不同的旋量结构选取 $\mathsf{s},\mathsf{s}'$ 相差一个元素 $\Delta\mathsf{s}\in H^1(M_3;\mathbb{Z}_2)$。现在，$\Delta\mathsf{s}$ 可以提升为具有整数周期的闭 1-形式 $\Delta\eta$。于是 $\eta(\mathsf{s})-\eta(\mathsf{s}')=\Delta\eta$ 可以有非平凡的周期，尽管只有这些周期模 2 的取值才有意义。

  \item 从旋量 Chern-Simons 项的定义容易看出它如何依赖旋量结构。若我们平移 $\epsilon\in H^1(M_3;\mathbb{Z}/2\mathbb{Z})$，则可以证明
  \begin{equation}
    \exp\left[2\pi i\frac12CS(A;\mathsf{s}+\epsilon)\right]
    =\exp\left[i\pi\int_{M_3}\epsilon\cup c_1(L)\right]\cdot
    \exp\left[2\pi i\frac12CS(A;\mathsf{s})\right]
    \tag{2.354}
    \label{eq:2.354}
  \end{equation}

  右端的前因子「只不过」是一个符号。为证明 \eqref{eq:2.354}，注意 $\eta(\mathsf{s})$ 的平移 $\Delta\eta$ 下降到 $\epsilon$，而 $\frac{F}{2\pi}$ 是 $c_1(L)$ 的代表元，其中 $L\to M_3$ 是 $M_3$ 上的主 $U(1)$ 丛。因此指数化作用量的改变（它只是一个因子）
  \begin{equation}
    \exp\left[i\pi\int_{M_3}\Delta\eta\wedge\frac{F}{2\pi}\right]
    \tag{2.355}
    \label{eq:2.355}
  \end{equation}
  可以简单地写成上同调中的杯积：
  \begin{equation}
    \exp\left[i\pi\int_{M_3}\Delta\eta\wedge\frac{F}{2\pi}\right]
    =\exp\left[i\pi\int_{M_3}\epsilon\cup c_1(L)\right]
    \tag{2.356}
    \label{eq:2.356}
  \end{equation}
\end{enumerate}

{\small\itshape ［原文标记（星号框）：评注：三维费米子路径积分的相位。APS 与 eta 不变量——解释指标定理的论证。］}
% [原文待办标记，原文 p.84]

\subsection{推广：环面规范群} % 原文: 2.6 Generalization: Torus Gauge Group

让我们把规范群推广为 $U(1)^d$。于是局部有 $d$ 个联络形式，局部写作 $A^I$，$I=1,\dots,d$，而一般的 Chern-Simons 作用量为
\begin{equation}
  S=\frac{1}{4\pi}\int_{M_3}K_{IJ}A^IdA^J
  \tag{2.357}
  \label{eq:2.357}
\end{equation}
其中 $K^{IJ}$ 是对称的整矩阵。

{\small\itshape ［原文标记：难道各个 $U(1)$ 因子不会通过单值性混在一起吗？］}
% [原文待办标记，原文 p.84；位于页边竖排]
% [Tutor Note: (2.357) 的积分核写的是 $K_{IJ}$（下标），而紧随其后的第 1--3 条与 (2.358) 写的是 $K^{IJ}$、$K^{II}$（上标）。在 $K$ 为对称整矩阵、$A^I$ 为 1-形式的约定下，这只是指标位置的排版不一致，数值上并无区别；但抄写时不要把它误读成 $K$ 与 $K^{-1}$ 的区别，也不要据此认为需要动度量来升降指标。]

\begin{enumerate}
  \item $K^{IJ}$ 决定一个整格 $\Lambda$。

  \item 对非旋量理论，对角矩阵元 $K^{II}$ 必须是偶数。所以我们谈的是偶整格。

  \item 对旋量理论，我们允许对角矩阵元 $K^{II}$ 为奇数。

  \item 旋量结构依赖：还需要为 $\Lambda$ 选取一个特征向量 $c_\Lambda$。于是 $c_\Lambda\otimes\widehat{w_2}$ 将是 $\Lambda\otimes H^2(Z_4;\mathbb{Z})$ 的特征向量，因此我们把作用量写成
  \begin{equation}
    \exp\left[2\pi i\int_{W_4}\frac12K_{IJ}\frac{F^I}{2\pi}\wedge
    \left(\frac{F^J}{2\pi}+c^J\otimes\widehat{w_2}\right)\right]
    \tag{2.358}
    \label{eq:2.358}
  \end{equation}

  \item Wilson 算符：特征标群 $\operatorname{Hom}(U(1)^d,U(1))\cong\mathbb{Z}^d$，故定义
  \begin{equation}
    W(\vec n,\gamma,f)=\exp\left[2\pi i\oint_\gamma n_IA^I\right]
    \tag{2.359}
    \label{eq:2.359}
  \end{equation}

  奇异的规范变换
  \begin{equation}
    A^I\to A^I+\ell^Id\phi
    \tag{2.360}
    \label{eq:2.360}
  \end{equation}
  其中 $\phi$ 是绕 $\gamma$ 的方位角，它诱导出等同
  \begin{equation}
    W(\vec n,\gamma,f)\sim W(\vec n+K\vec\ell,\gamma,f)
    \tag{2.361}
    \label{eq:2.361}
  \end{equation}

  {\small\itshape ［原文标记：需要讨论标架依赖性。］}
  % [原文待办标记，原文 p.85；位于页边竖排]

  \item 与外部规范场耦合：[9] 的「自旋-电荷关系」，另见 [55]。

  Belov-Moore 的量子等价陈述：
  \begin{enumerate}
    \item 带二次精化的判别式群（discriminant group）。
    \item Gauss-Milgram 公式。需要提升到 $c\bmod 24$。
    \item 一些非平凡的等价：$II^{1,1}$（Witten）以及 Niemeier 格。重访 Chetan Nayak 的文章。
  \end{enumerate}

  边界条件：
  \begin{enumerate}
    \item 非拓扑的：边缘态。
  \end{enumerate}

  当判别式群有平方根时：等价于有限群规范理论 [Maldacena-Moore-Seiberg; Banks-Seiberg]。

  1-形式对称性。

  时间反演。
\end{enumerate}

额外的外自同构：格 $\Lambda$ 的自同构作为相空间的外自同构而作用。它们如何作用在 Hilbert 空间上？

{\small\itshape ［原文标记：（星号分隔线）］}
% [原文标记，原文 p.85]

\subsubsection{对 FQHE 的应用} % 原文: 2.6.1 Applications To the FQHE

需要回顾多少 FQHE 的内容？

所谓「统计规范场」

对偶性。准粒子的密度。

\subsubsection{与外部阿贝尔规范场的耦合} % 原文: 2.6.2 Coupling To External Abelian Gauge Fields

Witten 的 $Sp(2n,\mathbb{Z})$ 作用。

\section{非阿贝尔规范场的 Chern-Simons 理论} % 原文: 3. Chern-Simons Theories For NonAbelian Gauge Fields

\begin{keytakeaways}
  \item 归一化：$\operatorname{tr}(x,y)=-\frac{1}{2h^\vee}\operatorname{Tr}_{\mathfrak{g}}\mathrm{Ad}(x)\mathrm{Ad}(y)$ 使紧致群上该双线性型正定、长余根的模方为 $+2$，从而 $\wp(A)=\frac{1}{8\pi^2}\operatorname{tr}F^2$ 的四维周期为整数，且当 $G$ 单连通时能取到 $1$。
  \item 局部代数：由 $\operatorname{tr}F(A+\alpha)^2-\operatorname{tr}F(A)^2=d\{\operatorname{tr}(2\alpha F+\alpha D_A\alpha+\frac23\alpha^3)\}$ 得到相对 Chern-Simons 形式 $CS(A+\alpha,A)$，并由此得到 $dCS(A_1,A_2)=\wp(A_1)-\wp(A_2)$。
  \item 整体障碍：主丛 $P\to M_3$ 非平凡时不存在整体定义的 $A$，故 $CS(A,0)$ 无意义；$\frac{1}{8\pi^2}\int_{M_3}\operatorname{tr}(AdA+\frac23A^3)$ 只在 $\mathbb{R}/\mathbb{Z}$ 中有定义，因此 $k$ 必须是整数。
  \item 量子化与反常：在大规范变换下作用量的移动由 Wess-Zumino 型项 $-\frac{1}{24\pi^2}\int\operatorname{tr}(g^{-1}dg)^3$ 的整数周期控制；作者在页边自注此处有「因子 2 的困惑」，说明其归一化并非对每个 $g$ 都显然给出任意整数。
\end{keytakeaways}

\subsection{一些 Chern-Weil 理论} % 原文: 3.1 Some Chern-Weil Theory

更多细节见附录 I。这里我们只总结若干要点。

我们最终会想讨论任意的紧致 Lie 群。

本节中，设 $G$ 是紧致、连通、单的 Lie 群，Lie 代数为 $\mathfrak{g}$。我们考虑维数 $\ge4$ 的流形 $M$ 上的主 $G$-丛 $P\to M$。我们在 $\mathfrak{g}$ 上引入一个 $\mathrm{Ad}$-不变双线性型
\begin{equation}
  \operatorname{tr}(x,y)=-\frac{1}{2h^\vee}\operatorname{Tr}_{\mathfrak{g}}\mathrm{Ad}(x)\mathrm{Ad}(y)
  \tag{3.1}
  \label{eq:3.1}
\end{equation}
其中 $h^\vee$ 是对偶 Coxeter 数。这是一个好的归一化，因为对紧致群它是正定的，而对单 Lie 代数，最长余根的模方为 $+2$。

{\small\itshape ［原文标记：核对！核对！］}
% [原文待办标记，原文 p.86；位于页边竖排]

{\small\itshape ［原文标记（星号框）：作为一项合理性检验，让我们考虑 $\mathfrak{g}=\mathfrak{su}(2)$ 的情形。］}
% [原文待办标记，原文 p.86]

另一个好处是：当 $G$ 还是单连通的时候，则
\begin{equation}
  \wp(A):=\frac{1}{8\pi^2}\operatorname{tr}F^2
  \tag{3.2}
  \label{eq:3.2}
\end{equation}
在所有四闭链上都具有整数周期，而且存在周期为 $1$ 的丛与四维流形。

现在，如附录 I 中详细解释的，对 $P$ 上的两个联络，我们有定义良好的形式 $CS(A_1,A_2)\in\Omega^3(M)$，使得
\begin{equation}
  dCS(A_1,A_2)=\wp(A_1)-\wp(A_2)
  \tag{3.3}
  \label{eq:3.3}
\end{equation}
若 $A_2=A$、$A_1=A+\alpha$，其中 $\alpha\in\Omega^1(M;\mathrm{ad}P)$ 是伴随丛 $\mathrm{ad}P$ 的整体良好定义的截面，则
\begin{equation}
  \operatorname{tr}F(A+\alpha)^2-\operatorname{tr}F(A)^2
  =d\left\{\operatorname{tr}\left(2\alpha F+\alpha D_A\alpha+\frac23\alpha^3\right)\right\}
  \tag{3.4}
  \label{eq:3.4}
\end{equation}
非正式地说，这就是
\begin{equation}
  \delta\left(\operatorname{tr}F^2\right)=2d\left(\operatorname{tr}\delta A\,F\right)
  \tag{3.5}
  \label{eq:3.5}
\end{equation}

{\small\itshape ［原文标记：把 $D_A\alpha$ 写明白。］}
% [原文待办标记，原文 p.86；位于页边竖排]

于是：
\begin{equation}
  CS(A+\alpha,A)=\frac{1}{8\pi^2}\operatorname{tr}\left(2\alpha F+\alpha D_A\alpha+\frac23\alpha^3\right)
  \tag{3.6}
  \label{eq:3.6}
\end{equation}

现在我们要写出一个只依赖单个规范场 $A$ 的 Chern-Simons 作用量。非正式地，我们想研究指数化的作用量：
\begin{equation}
  \exp\left[2\pi i\frac{k}{8\pi^2}\int_{M_3}\operatorname{tr}\left(AdA+\frac23A^3\right)\right]
  \tag{3.7}
  \label{eq:3.7}
\end{equation}

那么，为什么不直接考虑 $CS(A,0)$ 呢？若主 $G$ 丛 $P\to M_3$ 非平凡，则 $P$ 上根本没有 $A=0$ 这个联络这回事。只有当 $P$ 平凡时才能取平凡联络。一般地，3-形式 $\operatorname{tr}(AdA+\frac23A^3)$ 在 $M_3$ 上并不是定义良好的形式。尽管如此，积分表达式
\begin{equation}
  \frac{1}{8\pi^2}\int_{M_3}\operatorname{tr}\left(AdA+\frac23A^3\right)
  \tag{3.8}
  \label{eq:3.8}
\end{equation}
有一个定义良好的含义，只要我们只把它当作 $\mathbb{R}/\mathbb{Z}$ 的元素。因此，常数 $k$ 必须是整数。我们有三种不同的方式把 \eqref{eq:3.8} 定义为 $\mathbb{R}/\mathbb{Z}$ 的元素（三者彼此一致）：

\begin{enumerate}
  \item 我们把 $M_3$ 写成某个定向四维流形 $W_4$ 的边界，并把丛与联络延拓到 $W_4$ 上。（这种延拓的存在没有拓扑障碍。）然后我们定义
  \begin{equation}
    \int_{M_3}\frac{1}{8\pi^2}\operatorname{tr}\left(AdA+\frac23A^3\right)
    :=\int_{W_4}\wp(A)\ \in\mathbb{R}
    \tag{3.9}
    \label{eq:3.9}
  \end{equation}

  这里的问题在于：可能的延拓与 $W_4$ 的选取有很多。任意两种选取之间的差别将是：
  \begin{equation}
    \int_{W_4}\wp(A)-\int_{W'_4}\wp(A')=\int_{Z_4}\wp(A)\ \in\mathbb{Z}
    \tag{3.10}
    \label{eq:3.10}
  \end{equation}
  其中 $Z_4$ 是闭四维流形。由于 $\wp$ 的归一化使得我们可以得到任意整数（对单连通的 $G$），这些周期将是整数。因此，该定义只有作为 $\mathbb{R}/\mathbb{Z}$ 的元素才与延拓的选取无关。

  \item 虽然 $M_3$ 上没有整体定义的 $\mathfrak{g}$-值联络 1-形式，但在主 $G$-丛的全空间上有这样的联络。因此，在 $P$ 的全空间上存在定义良好的形式 $CS(A)$。\footnote{构造主丛全空间上整体定义的 Chern-Simons 形式，还有一个更优雅的办法：若 $\pi:P\to M$ 是主 $G$ 丛，则主 $G$ 丛 $\pi^*(P)\to P$ 可以典范地平凡化，因为 $s(p)=(p,p)$ 是一个整体截面。若 $A$ 是 $P\to M$ 上的联络，它将被拉回到 $\pi^*(P)\to P$ 上的联络，于是我们可以把 $CS(\pi^*(A),0)$ 定义为 $P$ 的全空间上整体定义良好的形式 $CS(A)$。}

  现在，如果我们能选取一个截面 $s:M\to P$，那么我们就可以拉回并积分 $\int s^*CS(A)$。这样做有两个问题：第一，若 $P\to M$ 是非平凡丛，则不存在连续截面。这不太严重，因为我们可以从 $M$ 中去掉一个测度为零的集合（在其上联络是光滑的），并在其余集上把丛平凡化。然而，所得的数将依赖于我们如何平凡化。若 $s_1,s_2$ 是两个不同的平凡化，它们相差 $s_1=s_2\cdot g$，其中 $g$ 是 $P$ 的自同构（规范变换）。在最简单（也最常见）的情形中，这意味着 $g:M\to G$ 是一个定义良好的映射。现在，当 $A$ 整体定义良好、从而 3-形式 $CS(A)$ 有意义时，容易验证如下规范变换性质：
  \begin{equation}
    CS(A^g)-CS(A)=CS(g^{-1}dg)+d\left(\frac{1}{8\pi^2}\operatorname{tr}Adg\,g^{-1}\right)
    \tag{3.11}
    \label{eq:3.11}
  \end{equation}
  以及
  \begin{equation}
    CS(g^{-1}dg)=-\frac{1}{24\pi^2}\operatorname{tr}\left(g^{-1}dg\right)^3
    \tag{3.12}
    \label{eq:3.12}
  \end{equation}
  因此 $\int s^*CS(A)$ 会改变 $-\frac{1}{24\pi^2}\int\operatorname{tr}(g^{-1}dg)^3$，而该形式在 $M$ 上的周期将是一个非零整数。于是，$\int_{M_3}CS(A)$ 有意义，但只作为 $\mathbb{R}/\mathbb{Z}$ 的元素。

  {\small\itshape ［原文标记：因子 2 的困惑。以我们的归一化来看，这似乎总是一个偶数！！！］}
  % [原文待办标记，原文 p.87；位于页边竖排]

\end{enumerate}
% 原文此处声称有「三种」把 (3.8) 定义为 $\mathbb{R}/\mathbb{Z}$ 中元素的方式，但只列出两种（第一种：写成四维流形的边界并延拓丛与联络；第二种：在主丛全空间上定义并在截面上积分），第三种在原稿中缺失。紧接着的段落转而讨论作用量自身的变分，并非第三种定义。

Chern-Simons 作用量本身如何变化？同样，利用附录中的公式，我们有
\begin{equation}
  \frac{d}{dt}CS(A+\alpha(t),A)
  =\frac{1}{8\pi^2}\left[2\operatorname{tr}\dot\alpha(t)F(A+\alpha(t))
  -d\left(\operatorname{Tr}\alpha(t)\dot\alpha(t)\right)\right]
  \tag{3.13}
  \label{eq:3.13}
\end{equation}
% 原 (3.13) 中第一个迹印作小写 $\operatorname{tr}$、第二个印作 $\operatorname{Tr}$，原稿混用了两种记号；此处照原样转写，应理解为同一个不变迹。
